\section*{The Propers: Detailed Commentary}

\subsection*{The 1962 ritual architecture: covenant brought to the altar}

The heading is precisely \latin{Pro sponsis}, ``for the spouses.'' General Rubrics 121 and 341--343 make it a white votive Mass of II class. Its own text directs that the \latin{Gloria} be said; no \latin{Credo} is said unless an occurring Sunday or octave would require it, although this formulary itself is never admitted on Sunday. The four controlling nuptial rubrics are compact but complete:

\begin{threestable}{Rubric}{Provision}{Consequence}
378 & The \latin{Pro sponsis} Mass---or at least its Collect when the Mass of the day intervenes---is permitted whenever nuptials are celebrated outside closed time, and also within closed time if the Ordinary, for a just cause, permits the solemn blessing. & Celebration of marriage and admission of this Mass / blessing are related questions, not one automatic permission.\\
379 & Beyond days already excluding II-class votives, the formulary is prohibited on every Sunday and whenever the blessing cannot be given under n.~381c. & A permitted marriage does not by itself make the votive formulary available.\\
380 & If the Mass is barred but the blessing allowed, the occurring Mass is used, the Nuptial Collect is added under one conclusion even where n.~343c would otherwise prevent it, and the blessing is given. If both are barred, Mass with blessing may follow on a suitable unimpeded day. & The proper can be impeded, commemorated, or transferred without treating the marriage itself as uncelebrated.\\
381 & The blessing is inseparable from Mass absent Apostolic indult (a); the Mass celebrant gives it even if another priest assisted at the marriage (b); it is omitted when either spouse is absent or, subject to the stated local custom, has already received it (c); All Souls and the Sacred Triduum exclude Mass, commemoration, and blessing (d). & These clauses govern the solemn nuptial blessing, not the validity of consent or every possible pastoral case.\\
\end{threestable}

Before this Mass, the governing book for contracting marriage is the 1952 typical edition of the \emph{Rituale Romanum}, Titulus VIII: cap.~I gives the preliminaries, cap.~II the \latin{Ritus celebrandi Matrimonii sacramentum}, and cap.~III the blessing outside Mass available only by Apostolic indult. The 1962 Missal's n.~381a itself sends the exceptional outside-Mass case to that third chapter. The Mass formulary begins at that handoff; banns, impediments, delegation, mixed marriage, local language and custom, and present authorization to use earlier books lie outside it.

The formulary supplies neither a proper Preface nor a referral to another nuptial Common. At the Preface, the seasonal rules in General Rubrics 484--490 and 493--494 control when applicable; otherwise the Common Preface is used under nn.~482 and 498. A later optional Nuptial Preface cannot be imported into this edition.

The nuptial action alters the ordinary sequence three times. After the Pater, before \latin{Libera nos}, the celebrant turns from the Epistle side toward the kneeling spouses and prays \latin{Propitiare} and the long \latin{Deus, qui potestate}. After receiving the Precious Blood, he communicates them. The complete ending is \latin{Ite, missa est}; \latin{Deus Abraham} over the spouses; admonition; holy-water sprinkling; \latin{Placeat tibi}; general blessing; and the usual Last Gospel beginning of St.~John (General Rubrics 507--509). The rite requires every link in this chain.

\begin{massbox}{What the architecture says without becoming Matrimony's form}
The covenant is brought inside the Eucharistic action: after the prayer Christ taught, before sacramental peace and Communion, and under the Church's final sending. This placement interprets married life as received blessing and ecclesial vocation. It does not make the prayers, priest, or Eucharist the efficient sacramental sign of Matrimony; consent has already established the bond.
\end{massbox}

\sectionguard
\subsection*{Tobias and Sarah: mercy in an exiled household\enspace\properrefs{Int.}}

The Introit is personal and ecclesial at once: \latin{Deus Israel coniungat vos}. It invokes Israel's covenant God over two persons, then remembers mercy shown to \latin{duobus unicis}, two only children. Within Tobit, Sarah's bridal chamber has been marked by death, accusation, and despair; Tobias travels under Raphael's concealed guardianship; both families pray, give thanks, and discover providence acting through vulnerable ordinary means. The Mass begins Christian marriage neither with romance as autonomous destiny nor with anxiety as final word, but with mercy and a petition for fuller praise.

The two recensions differ here. The Roman Introit follows the shorter Vulgate form at Tobit 7:15 and 8:19. Modern Catholic Bibles often follow the longer Greek recension, where the episodes align differently. The missal's citation and text are the proper's; the longer recension supplies narrative context, not a replacement quotation.

Bede's \emph{Allegorica interpretatio in Tobiam} treats Tobit VII--VIII at length (PL 91:930B--932A). He reads Raguel's careful consent as the nations' reception of the faith, Tobias and Sarah's union as Christ taking the Church from the Gentiles, and the bridal chamber's renunciation and Trinitarian confession as the sequence by which the adversary is expelled and believers are joined to Christ. He expounds the chapters that furnish the Introit directly, though not word for word through its two short-Vulgate clauses. Bede's ecclesial allegory therefore opens the literal exiled household toward Christ and the Church without replacing the text's human consent, danger, mercy, and praise.

Tertullian's description of a marriage approved by the Church, confirmed by oblation, sealed by blessing, and lived in one hope and service offers an early Christian analogue (\emph{Ad uxorem} II.8). It illuminates why Tobias's household prayer belongs at Mass, though it describes a second-century rite, not the 1962 formulary.

\begin{sourcecard}{Introit's complete movement}
\textbf{God of Israel} names covenant history. \textbf{Join you} assigns principal causality. \textbf{Be with you} makes union a common pilgrimage. \textbf{Mercy to two only children} remembers threatened lives received as gift. \textbf{Make them bless thee more fully} makes praise, not the couple itself, the final liturgical subject.
\end{sourcecard}

\sectionguard
\subsection*{The Collect and the ministers: office under blessing\enspace\properrefs{Coll.}}

\latin{Quod nostro ministratur officio} names a real public act of the Church. \latin{Tua benedictione potius impleatur} refuses to identify ritual competence with supernatural sufficiency. The comparative \latin{potius} places the ministerial action under the divine fulfillment for which it is an instrument.

In a Thomistic account, the spouses' present consent is the proximate instrumental cause of the marriage. Their words are sacramental form; their sensible acts stand in place of the material element familiar in other sacraments (\emph{ST} Suppl., q.~42, aa.~1--3; q.~45, a.~1). The priestly blessing is sacramental and ecclesially decisive for the rite's solemnity, yet not the essence of the bond. Scholastic schools differ in the precise vocabulary of matter and form; this account is the Thomistic school's, and not every detail of it is a separately defined dogma.

\begin{fourstable}{Cause / aspect}{Fulfillment here}{What it explains}{What it excludes}
Principal cause & God acting through Christ's Pasch and sacramental institution. & Why blessing can elevate nature and confer grace. & Magic, autonomous consent, or priestly ownership of the spouses' covenant.\\
Instrumental cause & Mutual, free, present consent of capable spouses. & Why the spouses minister the sacrament to one another. & A merely private contract or a bond imposed without consent.\\
Formal sign & Words outwardly expressing present consent. & How an interior act becomes sensible and ecclesially knowable. & Treating emotion, consummation, rings, or the nuptial Mass as the essential form.\\
Bond and grace & A stable relational unity and sacramental help for its acts and ends. & How the sacrament remains after the words and enables a shared vocation. & Fusion of persons, an indelible character, or guaranteed perseverance without cooperation.\\
\end{fourstable}

\sectionguard
\subsection*{Psalm 127: from fear and walking to Israel's peace\enspace\properrefs{Int., Grad., Tr., Comm.}}

The formulary distributes one psalm across the entire Mass. At entrance the blessed person fears the Lord and walks in his ways. At the Gradual the wife is fruitful vine and the children olive shoots around the table. In penitential seasons the Tract continues to Sion, Jerusalem's goods, generations, and Israel's peace. At Communion it compresses the same arc: God-fearing blessing, descendants, peace. What begins as moral motion becomes household abundance, then ecclesial and eschatological communion.

Augustine's direct exposition refuses to isolate these levels. Many walkers are one because Christ is one; the Church is both spouse and mother; children can signify those born to ecclesial life. Yet he also warns against a crude prosperity reading. Biological descendants are good, but their goodness lies in their being formed well; faithful childlessness is not divine rejection. Augustine keeps the psalm's literal blessing and opens it toward baptismal fruitfulness and heavenly Jerusalem.

\begin{fourstable}{Psalm image}{Literal integrity}{Ecclesial opening}{Pastoral consequence}
Fear and ways & Reverent obedience embodied through time. & The baptized walk as members of one Christ. & Fear is filial reverence, not terror or spousal intimidation.\\
Vine at the house & The spouse and household can bear real bodily and social fruit. & The Church bears children by word and sacrament. & A wife is a person and co-agent, never domestic property or a fertility instrument.\\
Olives at the table & Children share household nourishment and worship. & Anointing, endurance, and the ecclesial table can resonate typologically. & The text does not make every table image directly Eucharistic.\\
Sion and peace & Household blessing is received from Israel's sanctuary and people. & The domestic church is ordered to the whole Church's common good. & Family intimacy cannot become sectarian closure or excuse injustice.\\
\end{fourstable}

\sectionguard
\subsection*{The seasonal chants: one blessing under penance and Pasch\enspace\properrefs{Grad., All., Tr.}}

Outside Paschaltide the Gradual is followed by Psalm 19's Alleluia: help comes \latin{de sancto}, protection from Sion. Augustine interprets the sanctuary christologically as Christ's sanctified Body and links the royal prayer with sacrifice, Cross, Resurrection, and heavenly intercession (\emph{Exposition on Psalm 19/20}). The couple's request is thus ecclesial before it is sentimental.

After Septuagesima the Alleluia is omitted and Psalm 127 becomes the Tract. Its three verses slow the blessing into fear, persevering sight of Jerusalem's good, generations, and Israel's peace. In Paschaltide the Gradual is omitted; Psalm 19 remains as the first Alleluia and Psalm 133 adds blessing from Sion by the maker of heaven and earth. The risen Christ does not release marriage from creation: Paschal newness restores the Creator's purpose and carries it beyond death.

Augustine's direct exposition of Psalm 133 (titled 134), sections 1--3, gives this Paschal verse a demanding moral horizon: God's courts are charity's breadth rather than hatred's confinement; hands lifted by night bless in adversity; praise is contradicted when deeds curse. The Lord blesses from Sion because he makes many one: ``be one'' to receive the blessing. Augustine comments on the psalm, not the Nuptial Mass; applying his exposition to spousal works and unity is a reading of the verse's nuptial placement.

\sectionguard
\subsection*{The Gospel: the Creator's act against the logic of dismissal\enspace\properrefs{Gosp.}}

Matthew 19 begins with a test: may a man dismiss his wife \latin{quacumque ex causa}? Jesus does not arbitrate between looser and stricter catalogues of causes at their chosen level. He moves behind Deuteronomy's concession to Genesis: one Creator, male and female, leaving and cleaving, two becoming one flesh. The conclusion changes grammatical subject. Human beings do not manufacture an absolutely private arrangement; God joins, and man must not separate.

Chrysostom's \emph{Homily 62 on Matthew} makes the double force explicit. Separation is against nature because one flesh is cut, and against law because what God joined is divided. Yet the same discourse continues immediately to celibacy ``to whom it is given.'' Christ neither absolutizes every social custom surrounding marriage nor makes marriage the only complete Christian life.

St.~John Paul II's 1979 catecheses return to the same Gospel from a later historical horizon. He argues that Christ's appeal to the ``beginning'' is not nostalgia for an inaccessible innocence: historical men and women hear Genesis within the redemption of the body. He then reads one flesh as a chosen communion of persons in which the whole humanity of each spouse is mutually given. John Paul II thereby deepens the personal subjectivity of the bond while preserving the Creator's act and the Gospel's indissolubility; he does not replace Matthew's dispute with a modern anthropology.

Aquinas calls marriage a real relation: two persons are ordered to one generation and education and to one common domestic life (\emph{ST} Suppl., q.~44, a.~1). Their unity is therefore neither metaphor alone nor substantial fusion. It is a unity of order embodied in an exclusive relation whose justice and signification require permanence.

\begin{studybox}{What God has joined}
The formula \latin{Quod ergo Deus coniunxit} establishes divine authorship and human non-separation; it is not a command to conceal abuse, refuse protection, or confuse civil separation with a declaration that a valid bond never existed. Current law itself distinguishes separation with the bond remaining and recognizes grave mental or physical danger as a legitimate cause for leaving (CIC 1151--1155, especially 1153 \S1). Particular cases still call for competent pastoral or canonical judgment.
\end{studybox}

\sectionguard
\subsection*{Ephesians: sacrifice before analogy\enspace\properrefs{Ep.}}

The lection begins at Ephesians 5:22, but its immediate grammar is prepared by 5:2 and 5:18--21. Christ loves by handing himself over as sacrifice; the Spirit produces psalmody and thanksgiving; all submit to one another in the fear of Christ. The husband's comparison to Christ is therefore a judgment upon self-interest: love must sanctify, nourish, cherish, and give the self for the other's good.

Chrysostom's direct homily presses the consequence: no threats, insults, or violence; honor must be preferred to money; spouses should pray, worship, teach children, and make the house a little church. Some of his social assumptions belong to antiquity, but his Christological center directly resists domination. Aquinas likewise distinguishes governance for common good from rule for private utility and identifies three marital unities: charity, common life, and flesh.

John Paul II begins his 28 July 1982 Ephesians cycle by calling this passage the crown of Christ's words about the beginning, the heart, and resurrection. The analogy is therefore not a frozen allocation of social power. It belongs to the redemption of the body and is governed by mutual subjection in Christ's fear, Christ's saving self-gift, and the Church's holiness. Chrysostom supplies the ascetical household demands, Aquinas the ordered sacramental relation, and John Paul II the personal and redemptive horizon; their emphases converge without becoming one undifferentiated authority.

The wife's relation to the Church is also analogical, not reductive. The Church is responsive because Christ first creates, washes, and gives himself for her; she is a living body of members, not inert matter. A wife possesses equal baptismal dignity, intelligence, freedom, gifts, and direct relation to God. Ephesians' final command makes the husband love as himself and the wife reverence her husband; neither clause permits sin or erases reciprocal obligation.

\begin{sourcecard}{\latin{Sacramentum hoc magnum est}}
Paul identifies the ``great mystery'' ultimately with Christ and the Church. Marriage is great by signification because its one-flesh covenant points beyond itself; Catholic tradition also recognizes it as one of seven efficacious sacramental signs. The sign and the transcendent reality are distinct: the spouses do not reproduce the Pasch, and Christ's union with his Church does not depend upon their success.
\end{sourcecard}

\sectionguard
\subsection*{Offertory and Secret: all their times beneath the one oblation\enspace\properrefs{Off., Sec.}}

Psalm 30's speaker entrusts not merely a moment but \latin{tempora mea}, the succession of life, to God's hands. Augustine's direct exposition reads the psalm through Christ and the whole Church, especially Passion, deliverance, and election; his Psalter's \latin{sortes meae} differs lexically from the Roman Offertory, so his commentary supplies a Christological horizon rather than exact word study.

At precisely this chant, bread and wine are prepared. The spouses' embodied future---work, illness, delight, continence, children or childlessness, aging and death---can be offered analogically with the Church. The Secret calls the gift \latin{munus oblatum} for the sacred law of wedlock, then names God \latin{largitor} of the work and asks him to be \latin{dispositor}. Natural capacity, free promise, liturgical gift, providential ordering, and grace are coordinated without confusion.

Tertullian says ecclesial marriage is confirmed by oblation and lived by spouses who pray, fast, worship, and serve together. Aquinas sets the systematic order: the Eucharist is principal among the sacraments and their end; Matrimony serves the perfection of the multitude and touches the Eucharistic end by signifying Christ and the Church (\emph{ST} III, q.~65, aa.~2--3). The couple's self-offering is real moral participation, never another sacramental sacrifice alongside Christ's.

\sectionguard
\subsection*{After the Pater: author, helper, and preservation\enspace\properrefs{Nupt. I}}

The first nuptial prayer asks God to favor the Church's supplications and assist his own institutions by which he ordered propagation of the human race. Its final clause holds the balance: \latin{quod te auctore iungitur, te auxiliante servetur}. God authors the joining; God helps its preservation. The passive verbs hold the spouses within grace, while \latin{servetur} presupposes the active fidelity named by the final admonition.

Nature and grace are neither rival explanations nor parallel tracks. Generation belongs to living nature and serves continuation of the species; rational generation includes long education, household friendship, political society, and worship. In \emph{SCG} IV.78 Aquinas argues that Christian marriage orders offspring to the worship of God and the Church and, because sacraments effect what they signify, receives grace appropriate to that work. Grace heals desire and will, elevates friendship into charity, and orders temporal goods to beatitude.

\sectionguard
\subsection*{The order of ends: children, mutual aid, and beatitude\enspace\properrefs{Ep., Grad., Nupt. I, Nupt. II, Comm., Final}}

The first prayer after the Pater identifies propagation of the human race as an institution of providence. The Gradual and Communion place spouse, vine, olive shoots, table, generations, and Israel's peace in one psalmic arc; the long blessing asks for offspring and the final admonition commands diligent Catholic education. In the traditional Latin theology presupposed by this 1962 formulary, these texts manifest marriage's \textbf{primary natural and institutional end: the generation and education of children}. Generation is incomplete without the long work of forming a free person in truth, virtue, worship, and service.

The 1917 Code states that hierarchy at canon 1013 \S1, and Pius XI repeats it in \emph{Casti connubii} 17. The same tradition names as \textbf{secondary intrinsic ends} mutual aid / the good and mutual perfection of the spouses, together with the remedy or chaste ordering of concupiscence. The Epistle's sacrificial love, the prayers' author--helper grammar, love and peace, one bed, mutual fidelity, and conjugal chastity give those secondary ends real liturgical density. Their subordination does not make them optional and cannot justify using a spouse as a means; personal communion shelters, serves, and humanizes the primary end.

\Needspace*{7\baselineskip}
Current CIC 1055 \S1 and CCC 1601 ordinarily name the spouses' good and procreation / education together, while CCC 1652 calls children marriage's ``crowning glory'' and says the other ends are not diminished. The two vocabularies answer different questions: the older terminology formally ranks objective final causes; the current formulation protects the inseparable personal goods within the one covenant. Neither makes actual conception a validity test: CIC 1084 \S3 expressly states that sterility neither prohibits nor nullifies marriage. The covenant's kind and intention remain ordered to life even where providence does not grant children.

\begingroup
\scriptsize
\renewcommand{\arraystretch}{1.00}
\noindent\begin{tabular}{>{\bfseries\raggedright\arraybackslash}p{0.15\linewidth}>{\raggedright\arraybackslash}p{0.23\linewidth}>{\raggedright\arraybackslash}p{0.25\linewidth}>{\raggedright\arraybackslash}p{0.27\linewidth}}
\toprule
\textbf{Causal level} & \textbf{Matrimonial reality} & \textbf{Primary / secondary finality} & \textbf{Liturgical disclosure}\\
\midrule
Outward sign & Present consent under applicable Latin form & Consent establishes a covenant intrinsically ordered to its ends; it does not invent them & Collect and Secret distinguish ministry, divine gift, and providential ordering\\
\latin{Res et sacramentum} & The exclusive, enduring bond & Gives stable justice and common life for the reception and education of children and for mutual aid & God joins, preserves, connects, and guards in long peace\\
\latin{Res tantum} & Conjugal grace & Heals and elevates primary and secondary ends into parenthood, fidelity, chastity, mutual sanctification, and service & Christ's Bride, Eucharistic Communion, nuptial petitions, and admonition\\
Ultimate end & God's glory, persons, and charity, not the temporal institution as such & Children, spouses, and household are ordered beyond mortal succession to beatitude & Sion, Israel's peace, heavenly rest, kingdoms, and life without end\\
\bottomrule
\end{tabular}
\endgroup

Two uses of ``form'' meet at this Mass and remain distinct. Matrimony's sacramental form in the Latin analysis is the word / sign of consent; it effects a real bond and conjugal grace without converting either spouse's human substance. The Eucharistic consecratory form, by contrast, effects conversion of the whole substances of bread and wine into Christ's Body and Blood while the species remain. Thus covenant and transubstantiation are both real changes under divine causality, but changes of different metaphysical kinds.

\sectionguard
\subsection*{The long blessing: salvation history gathered over the spouses\enspace\properrefs{Nupt. II}}

The prayer's threefold invocation of \latin{Deus} builds a miniature economy of salvation.

\begin{threestable}{Movement}{Appointed language}{Doctrinal horizon}
Creation & God makes all from nothing, orders the beginnings of the universe, makes man in his image, and forms inseparable help from human flesh. & Contingent being, providential order, equal creaturehood, sexual difference, social unity, and one-flesh indissolubility.\\
Mystery & God consecrates the conjugal joining to prefigure the sacrament of Christ and the Church in the nuptial covenant. & Natural sign is elevated into an efficacious ecclesial sign under Christ's Pasch.\\
Blessing through judgment & The primordial social blessing alone is said to survive both original sin's penalty and the Flood's sentence. & Sin wounds but does not annihilate the created good; covenant history preserves and disciplines it.\\
Petition for the bride & Love and peace, fidelity and chastity, Rachel, Rebecca and Sarah, one bed, heavenly teaching, offspring, innocence. & Primary procreative / educative and secondary unitive / remedial ends are blessed together without losing their order.\\
Common final end & Both see generations, reach desired old age, the rest of the blessed, and heavenly kingdoms. & Temporal flourishing remains penultimate and is taken into eternal beatitude.\\
\end{threestable}

The blessing speaks at greatest length over the bride because of the inherited Roman form; that ritual asymmetry cannot erase the mutual duties stated by Scripture, consent, and the concluding admonition. Rachel, Rebecca, and Sarah are not interchangeable ideals of temperament. They locate this marriage inside the difficult providence of covenant households: fertility and barrenness, migration, conflict, promise, and fidelity carried through weakness.

\Needspace*{3\baselineskip}
The blessing joins creation to redemption across sin and Flood. Marriage is neither mere remedy nor self-sufficient: its wounded created good endures under covenant and Christ consecrates it.

\sectionguard
\subsection*{Communion and Postcommunion: table, bond, and peace\enspace\properrefs{Comm., Postcomm.}}

The missal directs the celebrant, after receiving the Precious Blood, to communicate the spouses and continue Mass. Psalm 127 then calls blessed every person who fears the Lord, invokes descendants, and ends with peace upon Israel. Communion locates the couple inside the one Body; it does not canonize the household as self-contained.

The Postcommunion's metaphysics are relational and providential. God accompanies \latin{instituta providentiae tuae} with gracious favor; he \latin{connectis} the spouses in \latin{legitima societate}; he is asked to \latin{custodias} them in \latin{longaeva pace}. Marriage is neither a momentary intention nor a new substance. It is a stable relation and common life whose moral acts remain free, whose peace requires virtue, and whose sacramental grace supplies help without eliminating contingency.

The Eucharist strengthens charity and ecclesial unity, the virtues by which marital works become meritorious. But sacramental reception does not mechanically repair injustice without repentance. The spouses remain responsible to seek reconciliation, protection, counsel, penance, and conversion as needed.

\sectionguard
\subsection*{The Abrahamic blessing and final admonition\enspace\properrefs{Final}}

After the dismissal the celebrant invokes the God of Abraham, Isaac, and Jacob. Three and four generations represent temporal continuity; \latin{et postea} decisively turns beyond it to eternal life without end. The serious admonition then translates sacramental theology into a rule of life: remain in God's fear, love one another, preserve mutual fidelity and conjugal chastity, and educate children diligently in the Catholic religion. Holy-water aspersion, \latin{Placeat}, the general blessing, and St.~John's prologue return this particular rule to the Church's ordinary conclusion of Mass.

The naming makes Christian marriage heir to concrete covenant history rather than a timeless archetype. Education completes the primary end named by propagation; \latin{mutuam fidem} and chastity articulate intrinsic secondary ends that protect persons and household. Formation is not possession of children but service to their natural and supernatural good, while the completed end sequence refuses to stop with temporal descendants: Christ is the helper toward life without end.

\begin{massbox}{Cross-sacramental economy}
\textbf{Baptism} forms the washed Bride and makes children adopted heirs. \textbf{Confirmation} strengthens household witness. \textbf{Eucharist} nourishes charity and is the rite's liturgical end. \textbf{Penance} restores truth and forgiveness. \textbf{Anointing} accompanies weakness and death. \textbf{Orders} serves the ecclesial Body. \textbf{Matrimony} serves it primarily through generation and education, secondarily through mutual aid and chaste communion. All are ordered through Christ to final beatitude.
\end{massbox}

\sectionguard
\subsection*{Marriage, continence, and the teleological end of all things}

Marriage and consecrated continence name two modes of bodily discipleship. Augustine calls both goods while distinguishing their signs: marriage signifies the enduring unity of Christ and the Church through fidelity in time; continence anticipates directly the resurrection where no one marries. The higher eschatological sign does not make marriage evil, and marriage's sacramentality does not make continence defective.

Aristotle sees generation as a mortal creature's participation, as far as possible, in continuity and even in the divine; Christian revelation both receives and judges that natural teleology. No succession of descendants defeats death. The final blessing therefore moves from generations to eternal life, and the long blessing from old age to heavenly kingdoms. Creation's generative good reaches its proportionate natural ends, while grace carries persons beyond species-continuity into vision of God.

Revelation 19 and 21 give the final horizon: the Lamb's marriage supper and the New Jerusalem descending as Bride. The Church's nuptial future is the consummation of creation, covenant, Incarnation, Cross, Resurrection, sacrament, household holiness, and cosmic renewal. Earthly marriage serves that end precisely by not claiming to be it.
